\section{\secdivalg}
\label{sec_div_alg}
 
In this section, we continue our discussion of integer division.  We introduce vocabulary to describe the result of dividing one integer by another when the answer is not an integer.  Next, we revisit polynomials and use polynomials to prove facts about even, odd, and divisors using cases.

%%%%%%%%%%%%%%%%%%%%%%%%%%%%%%%%%%%%%%%
\newcommand{\subdivalg}{The Division Algorithm, \fn{div}, and \fn{mod}}
\subsection{\subdivalg}
\label{sub_division_algorithm}

As we saw in \Cref{sec_even_odd_divides}, sometimes when we divide one integer by another integer we get an integer answer, such as $$40 \div 5 = \frac{40}{5} = 8.$$ We said that 5 divides 40 and wrote $5 \mid 40$. Other times, we get a fraction that is not an integer, such as $$23 \div 4= \frac{23}{4} = 5.75.$$ We said that 4 does not divide 23 and wrote $4 \nmid 23$. How can we think of $23 \div 4$ in terms of integers?  

\begin{exam}[Cookies and Children]
\label{exam_cookies_children}
    Use the idea of giving 23 cookies to four children to describe $23 \div 4$ in terms of (only) integers.
    \begin{soln}
        Imagine that I have 23 cookies and four children.  I want to give each child as many cookies as I can, but, to be fair, each child should have the same number of cookies as the other children.  What should I do?  I can give each child five cookies. That accounts for $5 \times 4 = 20$ of the cookies, but I cannot evenly distribute the remaining $23-20=3$ cookies without breaking them, which I do not want to do.  I guess those 3 cookies are for me to eat.   We might say $23 \div 4$ is 5 cookies/child with 3 cookies remaining for me.  We can write the equation $$23 = 5 \times 4 + 3.$$
    \end{soln}
\end{exam}

We formalize the idea from \Cref{exam_cookies_children} with the following theorem.  

\begin{thm}[The Division Algorithm]
\label{thm_div_alg}
Given integers $a$ and $n$ with $n \ge 1$, when we divide $a$ by $n$ we get $$a = qn + r$$ for some integers $q$ and $r$ such that $0 \le r < n$.  Moreover, $q$ and $r$ are unique.

\begin{proof}
    Let $a$ and $n$ be integers with $n \ge 1$.  We prove existence in the case where $a \ge 0$. % SU:  I'm not proving when $a$ is negative or the uniqueness.
    
    Since $a$ is nonnegative and $n \ge 1$, we can imagine giving $a$ cookies to $n$ children. First, I check to see if I have enough cookies for each child to get at least one, and if so, then I hand one cookie to each child.  Next, I repeat that step with the remaining cookies as many times as possible.  In the end, let's say that each child gets $q$ cookies.  In total, I gave $q \times n=qn$ cookies to the children, so there are $$a-qn=r$$ cookies leftover for me.  Adding $qn$ to each side of the equation we get $$a = qn+r.$$  Because I checked there were enough cookies for each child to get one, I never run out of cookies, and so $r \ge 0$.  On the other hand, any time I still have enough cookies to give each child another cookie, I would do that.  Therefore, the final number of remaining cookies must be strictly less than the number of children and so $r <n$.    
\end{proof}
\end{thm}

The integers $q$ and $r$ in the Division Algorithm (\Cref{thm_div_alg}) have special names.

\begin{defn}[Quotient, Remainder, \fn{div}, and \fn{mod}]
\label{defn_quo_rem_div_mod}

Let $a$ and $n$ be integers with $n \ge 1$ and let $q$ and $r$ be integers such that $a = qn+r$ and $0 \le r < n$.
\begin{apart}
    \item The integer $q$ is the \vocab{quotient} of $a \div n$. When $a \ge 0$, it is the number of cookies per child in the story. For example, the quotient when $23 \div 4$ is five.  
    \item We write $a$ \fn{div} $n =q$.  For example,  $23$ \fn{div} $4=5$.  
    \item The integer $r$ is the \vocab{remainder} of $a \div n$.  When $a \ge 0$, it is the number of cookies for me.  For example, the remainder when $23 \div 4$ is three.  
    \item We write $a$ \fn{mod} $n=r$.  For example, $23$ \fn{mod} $4=3$. In many computer programming languages, \fn{mod} is denoted using the percent symbol (\%).  For example, we would write $23 \% 4 = 3.$
    \item We say $a \div n$ is $q$ \vocab{with remainder} $r$. For example, $23 \div 4$ is five with remainder three.
    \item In \vocab{long division} the quotient is typically written at the top of the bar and the remainder is the final number at the bottom, as in \Cref{fig_longdivision_23div4}.
       \begin{figure}[hbtp]
           \centering
           \intlongdivision{23}{4} 
           \caption{Long division calculation for $23 \div 4$.}
           \label{fig_longdivision_23div4}
       \end{figure}
    \end{apart}
\end{defn}

Let's see these definitions in action.

\begin{exam}[Quotient, Remainder, \fn{div}, and \fn{mod}]
\label{exam_quo_rem_div_mod}
On this problem we do not use computational technology (except to check).
\begin{apart}
    \item Find the quotient and the remainder when $100 \div 9$.  Use your answer to calculate $100$ \fn{div} $9$ and $100$ \fn{mod} $9$.
        \begin{soln}
            Starting with 100 cookies, we recognize that $9 \times 11 = 99$ and so we can hand each child 11 cookies, leaving us with $100 - (11\times 9) = 100-99 = 1$ cookie.

            If we were using computational technology, we could calculate $100 \div 9 = 11.1111...$ and again see that 11 is the largest number of cookies we can give each child, as there are not enough cookies to give each child 12 cookies.
            
            If we did a long division instead, we would get the same answer, as shown in \Cref{fig_longdivision_23div4}.
            \begin{figure} [hbtp]
                \centering
                \intlongdivision{100}{9} 
                \caption{Long division calculation for $100 \div 9$.}
                \label{fig_longdivision_100div9}
            \end{figure}
            Notice that in the long division we calculate the remainder in two steps. First, we give each child 10 cookies for a total of $10 \times 9 = 90$ cookies (although 0 is not written), so the remainder is $100-90 = 10$.  Next, we give each child 1 cookie for a total of $1 \times 9 = 9$ cookies, so the final remainder is $10-9=1$.  
            
            Using any way of thinking of this division, we get $$ 100 = 11(9) +1,$$ so the quotient is 11 and the remainder is 1.  Thus, $100$ \fn{div} $9=11$ and $100$ \fn{mod} $9=1$.
        \end{soln}
    \item Find the quotient and the remainder when $35 \div 5$.  Use your answer to calculate $35$ \fn{div} $5$ and $35$ \fn{mod} $5$.
        \begin{soln}
            Write $35 = 7(5)$ so the quotient is seven.  To find the remainder, think of the equation as $35 = 7(5)+0$, so the remainder is zero.  Thus, $35$ \fn{div} $5=7$ and $35$ \fn{mod} $5=0$
        \end{soln}     
    \item Use the \fn{mod} notation (or \%) to describe whether an integer is even or odd.
        \begin{soln}
            An integer $n$ is even if $n=2k$ for some integer $k$.  When we divide, we get $n \div 2 = k$ and the remainder is zero. Thus, $n$ is even exactly when $n$ \fn{mod} $2=0$ or, equivalently, when $n\%2=0$.

            An integer $n$ is odd if $n=2k+1$ for some integer $k$.  When we divide, we get $n \div 2 = k$, and the remainder is one. Thus, $n$ is even exactly when $n$ \fn{mod} $2=1$ or, equivalently, when $n\%2=1$.
        \end{soln}
\end{apart}
\end{exam}

It is your turn to practice.

\begin{act}[Practice with Div and Mod]
\label{act_practice_div_mod}
In this activity, do all calculations by hand.  
    \begin{actpart}
        \item Find the quotient and the remainder when $73 \div 7$.  State your answer using \fn{div} and \fn{mod}.
        \item Find the quotient and the remainder when $66 \div 6$.  State your answer using \fn{div} and \fn{mod}.  
        \item Find the quotient and the remainder when $3 \div 10$.  State your answer using \fn{div} and \fn{mod}.  
        \item \label{part_mod10} Calculate 
        $27$ \fn{mod} 10, 115 \fn{mod} 10, and \text{123,456,789} \fn{mod} 10.  Hint: Subtract multiples of 10.  
        \item What do you notice in (\ref{part_mod10})?  State your answer as a conjecture.
        \item For a nonnegative integer $a$, what are the possible values of $a$ \fn{mod} $10$?  State your answer as a conjecture.
        \item For a nonnegative integer $a$ and an integer $n \ge 1$, what are the possible values of $a$ \fn{mod} $n$? State your answer as a conjecture.
    \end{actpart}
\end{act}

%%%%%%%%%%%%%%%%%%%%%%%% %%%%%%%%%%
\newcommand{\subpolys}{Algebra: Polynomials}
\subsection{\subpolys}
\label{sub_polys}

The Division Algorithm (\Cref{thm_div_alg}) gives us expressions such as $2k+1$ for an odd integer or $3q+2$ for an integer where that integer $\fn{mod}~ 3 = 2$. We might square an even integer to get an expression of the form $4n^2$. These expressions are polynomials. In this section, we take a look at polynomials, including how to simplify products and powers. We start with the definition. 

\begin{defn}[Polynomials]
\label{defn_polynomials}
~
    \begin{apart}
        
        \item For an integer $n$, a \textbf{polynomial} in $n$ is a sum of constant multiples of powers of $n$.  For example, $2m^4+5m^2-m+7$ is a polynomial in $m$.  To see why, we can rewrite this polynomial to show the powers of $n$.
        $$2m^4+5m^2-m+7=2(m^4)+5(m^2)+(-1)(m^1)+7(m^0).$$
        
        \item In a polynomial, the integer in front of a power of $n$ is the \vocab{coefficient} of that power of $n$.  For example, in the polynomial $2m^3+5m^2-m+7$, the coefficient of $m^4$ is 2, the coefficient of $m^2$ is 5, the coefficient of $m$ is $-1$, and the coefficient of $m^3$ is 0.
        
        \item In a polynomial, the coefficient of $m^0$ is the \vocab{constant term}. For example, the constant term of $2m^3+5m^2-m+7$ is 7.

        \item We add polynomials by \textbf{collecting like terms}.  For example $$(k^3 + 2k -7) + (4k^3 + k^2 + 2) = (k^3+4k^3) + (k^2) + (2k) + (-7 +2) = 5k^3+k^2+2k-5$$
        
        \item To multiply or \vocab{expand} a product of polynomials, we use the distributive rule.  See, for example, \Cref{exam_mult_polys}.
    \end{apart}
\end{defn}

We use the distributive rule (\Cref{defn_integer_algebra_order_operations})
to expand polynomials.

\begin{exam}[Multiplying polynomials]
\label{exam_mult_polys}  
~
    \begin{apart}
        \item Expand $n^2(2n+1)$.
            \begin{soln}
                By the distributive rule, 
                $$n^2(2n+1) = n^2(2n) + n^2(1) = 2n^3+n^2.$$
            \end{soln}
        \item Expand $(n+2)(n+3)$.
            \begin{soln}
                By the distributive rule (repeatedly),
                \begin{center}
                    \begin{tabular}{lll}
                    $(n+3)(n+2)$ &=& $n(n+2) + 3(n+2)$ \\ 
                    &=& $\underbrace{(n)(n)}_{\text{First}} 
                    +\underbrace{(n)(2)}_{\text{Outside}} 
                    + \underbrace{(3)(n)}_{\text{Inside}} + \underbrace{(3)(2)}_{\text{Last}}$ \\ 
                    &=& $n^2 + 2n + 3n + 6$ \\
                    &=& $n^2+5n + 6$           
                    \end{tabular}
                \end{center}
                The words under each product indicate the original location of the two integers that are multiplied. This method is \vocab{First, Outside, Inside, Last} or \vocab{FOIL}. 
            \end{soln}  
    \end{apart}
\end{exam}

There is a nice set of visualizations\footnote{This lesson is from the College Algebra Foundations course by Lumen Learning retrieved from \url{https://tinyurl.com/vizpolymult}.} of multiplying polynomials in $x$ as areas at \url{https://tinyurl.com/vizpolymult}.  Try your hand at multiplying polynomials with discrete variables.

\begin{act} [Multiplying polynomials]
\label{act_mult_polys}
    Expand and simplify each expression.
    \begin{apart}
        \item $2n(3n^2+2n-1)$
        \item $(k-2)(k+1)$
        \item $(m+1)^2$
        \item $(2q+1)^2$
        \item $(n+3)(n-3)$  
        \item $(n+b)(n-b)$ for any integer $b$
    \end{apart}    
\end{act}

%%%%%%%%%%%%%%%%%%%%%%%%%%%%%%%%%%%%%%%
\newcommand{\subpffcases}{Proof Format: Proof by Cases}
\subsection{\subpffcases}
\label{sub_pff_cases}

We can use the Division Algorithm (\Cref{thm_div_alg}) to prove that every integer is even or odd, but never both.

\begin{exam}[Every integer is even or odd]
\label{exam_every_integer_even_or_odd}
    Use the Division Algorithm (\Cref{thm_div_alg}) to prove that every integer is even or odd.
    \begin{soln}
        \emph{Proof}  Let $a$ be an integer.  By the Division Algorithm (\Cref{thm_div_alg}) with $n=2$, there exist unique integers $q$ and $r$ such that $a=qn+r$ and $0 \le r < 2$.  Note that $r=0$ or $r=1$.  Since $r$ is unique, we are in exactly one of these two cases.
        Case 1: Assume $r=0$.  Then $a = 2q+0 = 2q$ is even, by the definition of even.
        Case 2: Assume $r=1$. Then $a=2q+1$ is odd, by the definition of odd.  In either case $a$ is even or $a$ is odd. 
    \end{soln}
\end{exam}

As in \Cref{exam_every_integer_even_or_odd}, the Division Algorithm (\Cref{thm_div_alg}) can be used to break a proof into cases, similar to how we break counting arguments into cases in \Cref{sec_counting_steps_cases}.  Here is the proof format.

\begin{pff} [Proof by Cases]
\label{pff_cases} 
Prove: $P$
    \begin{proof}
        (Explain why there are $n$ cases.) We consider $n$ cases.
  
        Case 1: Assume \ldots (Explain why $P$ is true in this case.) 
        
        Case 2: Assume \ldots (Explain why $P$ is true in this case.) 
        
        \ldots

        Case $n$: Assume \ldots (Explain why $P$ is true in this case.) 
        
        In (either case/all cases) we have $P$. 
\end{proof}
\end{pff}

In a proof by cases, it is important that the cases cover all of the possibilities.  Sometimes we need to explain that up front.  The phrase ``We consider $n$ cases'' that alerts the reader that we are using a proof by cases.  The assumption for each case ``Assume that \ldots''.  is stated as a complete sentence (with a period at the end) before any additional argument is presented.  After we finish the cases, we summarize using the phrase ``in either case''.  If there are more than two cases, we use ``in all cases''.

\begin{act}[Proof by Cases: $n^2 + n$ is even]
\label{act_proof_cases_nsqplusn_even} 
~
    \begin{actpart}
        \item \label{part_proof_cases_read} Copy the following example of a proof by cases.
        
            Prove that $n^2+n$ is even for any integer $n$. 
            
            \begin{proof}               
                Let $n$ be an integer.  By \Cref{exam_every_integer_even_or_odd}, $n$ is even or $n$ is odd.  We consider two cases.
            
               Case 1: Assume $n$ is even.  By the definition of even, $n = 2k$ for some integer $k$.  Thus $$n^2+n=(2k)^2+(2k) = 4k^2+2k = 2(2k^2+k),$$ so $n^2+n$ is even. 
           
               Case 2: Assume $n$ is odd.  By the definition of odd, $n=2k+1$ for some integer $k$.  Thus $$n^2+n= (2k+1)^2 +(2k+1) = 4k^2 + 4k + 1 + 2k + 1 = 4k^2 + 6k + 2 = 2(2k^2 + 3k + 1),$$ so $n^2+n$ is even. 

               In either case, $n^2+n$ is even.         
            \end{proof}
        
        \item Edit the proof in (\ref{part_proof_cases_read}) to prove that $n^2 - n$ is even for any integer $n$.         
    \end{actpart}   
\end{act}

We give another example of proof by cases.

\begin{exam}[Using proof by cases to prove multiple of three.]
\label{exam_product_3consec_mult3}
    Use a proof by cases to prove 
        $$3 \mid n(n+1)(n+2)$$ for any integer $n$. 
        \begin{soln}
            Let $n$ be an integer.  By the Division Algorithm (\Cref{thm_div_alg}), we can write $$n=3q+r$$ for some integers $q$ and $r$ with $0 \le r < 3$.  The only possible values of $r$ are $r=0$, $r=1$, and $r=2$. We consider three cases.

            Case 1: Assume $r=0$.  In this case, $n=3q$ and so
                $$n(n+1)(n+2)= (3q)(3q+1)(3q+2) = 3[q(3q+1)(3q+2)] = 3k$$
            where $k = q(3q+1)(3q+2)$. Thus, $3\mid n(n+1)(n+2)$.

            Case 2: Assume $r=1$.  In this case, $n=3q+1$ and so
                $$n(n+1)(n+2)= (3q+1)(3q+2)(3q+3) = (3q+1)(3q+2)3(q+1) =3[(3q+1)(3q+2)(q+1)] = 3k$$
            where $k = (3q+1)(3q+2)(q+1)$. Thus, $3\mid n(n+1)(n+2)$.

            Case 3: Assume $r=2$.  In this case, $n=3q+2$ and so
                $$n(n+1)(n+2)= (3q+2)(3q+3)(3q+4) = (3q+2)3(q+1)(3q+4) = 3[(3q+2)(q+1)(3q+4)] = 3k$$
            where $k = (3q+2)(q+1)(3q+4)$. Thus, $3\mid n(n+1)(n+2)$.

            In all cases, $3\mid n(n+1)(n+2)$.
        \end{soln}
\end{exam}

%%%%%%%%%%%%%%%%%%%%%%%%%%%%%%%%%%%%%%%%%%

\subsection{Exercises}

%%%%%%%%%%%%%%%%%%%%%%%%%%%%%%%%%%%%%%%%%%

\subsubsection*{Exercises for \Cref{sub_division_algorithm} \subdivalg}

\begin{exer}[\exerP] % finding quotients and remainders]
\label{exer_find_quotient_remainder}
Do not use technology or other resources.
    \begin{exerpart}
        \item Find the quotient and the remainder when $100 \div 3$.  State your answer using \fn{div} and \fn{mod}.
        \item Calculate 20 \fn{div} 7 and 20 \fn{mod} 7.
        \item Explain your answers to (b) using the context of children and cookies as in \Cref{exam_cookies_children}.  
    \end{exerpart}
\end{exer}

\begin{exer}[\exerP] % names for mod 2]
\label{exer_mod2_even_odd}
~
    \begin{exerpart}
        \item Evaluate each of the following integers: 3 \fn{mod} 2, 4 \fn{mod} 2, 5 \fn{mod} 2, 6 \fn{mod} 2, 28 \fn{mod} 2, 125 \fn{mod} 2.
        \item Which integers have $n$ \fn{mod} $2 = 0$? 
        \item Which integers have $n$ \fn{mod} $2 = 1$? 
    \end{exerpart}
\end{exer}

\begin{exer}[\exerU] % Examples mod
\label{exer_examples_mod}
    On this problem, no justification is required.
    \begin{exerpart}
        \item Give an example of an integer $n$ with $100 \le n \le 200$ and $n$ \fn{mod} 2 = 1.
        \item Give an example of an integer $n$ with $100 \le n \le 200$ and $n$ \fn{mod} 3 = 1.
        \item Give an example of an integer $n$ with $100 \le n \le 200$ and $n$ \fn{mod} 7 = 1.
        \item Give an example of an integer $n$ with $100 \le n \le 200$ and $n$ \fn{mod} 10 = 1.
    \end{exerpart}
\end{exer}

\begin{exer}[\exerU] % Mod 10 practice]
\label{exer_mod10_practice}
~
    \begin{exerpart}
        \item Evaluate each of the following: 21 \fn{mod} 10, 37 \fn{mod} 10, 43 \fn{mod} 10, 59 \fn{mod} 10, and 60 \fn{mod} 10.
        \item Give an example of positive integers $a$ and $b$ where $2a \fn{ mod } 10 = 2b \fn{ mod } 10$ but $a \fn{ mod } 10 \neq b \fn{ mod } 10$.
        \item Give an example of integers $m$ and $n$ where $n \fn{ mod } 10 \neq 0$ and $m \fn{ mod } 10 \neq 0$, but $nm \fn{ mod } 10 = 0$.
    \end{exerpart}
\end{exer}

\begin{exer}[\exerU] % mod $n$]
\label{exer_modn}  
    Consider an integer $n \ge 1$.
    \begin{exerpart}
        \item What are the values of $0$ \fn{div} $n$ and $0$ \fn{mod} $n$?  Explain.
        \item What are the values of $n$ \fn{div} $n$ and $n$ \fn{mod} $n$?  Explain.
        \item If $a$ is an integer with $0 \le a < n$, what can you say about $a$ \fn{div} $n$ and $a$ \fn{mod} $n$?  Explain.
    \end{exerpart}
\end{exer}

\begin{exer}[\exerU] % using the Division Algorithm to generate cases]
\label{exer_cases_mod3_mod10}
~
    \begin{exerpart}
        \item \label{part_div_alg_mod3} Use the Division Algorithm (\Cref{thm_div_alg}) to explain how any integer can be written in the form $3q$, $3q+1$, or $3q+2$.  Hint: We used this fact in \Cref{exam_product_3consec_mult3}.
        \item \label{part_div_alg_mod10} Use the Division Algorithm (\Cref{thm_div_alg}) to explain how any integer $n$ can be written in the form $n=10q+s$ where $s$ is a digit, namely the last digit of $n$.  
    \end{exerpart}
\end{exer}

\begin{exer}[\exerR] % The Division Algorithm, div, and mod
\label{exer_dyk_div_alg}
    Do you know \ldots
        \begin{exerpart}
            \item What the Division Algorithm says?
            \item What restrictions are on the remainder $r$ when we calculate $n \div a$?
            \item How to calculate \fn{div} and \fn{mod} using long division?
        \end{exerpart}
\end{exer}

\begin{exer}[\exerE] % calculating \fn{div} and \fn{mod} from division]
\label{exer_calculate_div_mod_division}
~
    \begin{exerpart}
        \item Calculate 247 \fn{div} 18 and 247 \fn{mod} 18.
        \item Note that $247 \div 18= 13.7222\ldots$.  Explain how we could calculate 247 \fn{div} 18 using this equation.
        \item In general, if we know the decimal number $n \div a$, how can we find $a$ \fn{div} $n$? You may describe your answer in words.
        \item Again, note that $247 \div 18= 13.7222\ldots$.  Explain how we could calculate 247 \fn{mod} 18 using this equation. 
        \item In general, if we know the decimal number $n \div a$ and we know $a$ \fn{div} $n$, how can we find $a$ \fn{mod} $n$. You may describe your answer in words.
    \end{exerpart}
\end{exer}

\begin{exer}[\exerE] % ISBN numbers
\label{exer_isbn}
    The International Standard Book Number (ISBN) is a 13-digit integer used to identify a book.  The last digit is calculated from the first 12 digits and is used to check whether a 13-digit number is a legitimate ISBN.  This \textbf{check digit} is useful to avoid copying errors. Here is how.  For 
        $$\text{ISBN } = abcdefghijkmx$$
    the 13th check digit $x$ is chosen so that 
        $$(a + 3b + c + 3d + e + 3f + g + 3h + i + 3j + k + 3m + x) \fn{ mod } 10 = 0$$
            \begin{exerpart}
                \item Confirm that 9780691151649 is a valid ISBN.  Note that $a=9$, $b=7$, \ldots, $x=9$.
                \item Calculate the check digit for the ISBN beginning 123450123450.
            \end{exerpart}
\end{exer}

%%%%%%%%%%%%%%%%%%%%%%%%

\subsubsection*{Exercises for \Cref{sub_polys} \subpolys}

\begin{exer}[\exerP] % polynomial vocabulary]
~
    \begin{exerpart}
        \item What is the coefficient of $n^2$ in the polynomial $5n^3-n^2+3n+8$?
        \item What is the constant term of the polynomial $5n^3-n^2+3n+8$?
        \item What is the coefficient of $k$ in the polynomial $5k^2+k$?
        \item What is the constant term of the polynomial $5k^2+k$?
    \end{exerpart}
\end{exer}

\begin{exer} [\exerP] % expanding polynomials]
\label{exer_expanding_polys}
    Expand and simplify.
    \begin{exerpart}
        \item $n(2n^2+1)$
        \item $(k+1)(k+3)$
        \item $(m-1)(m+1)$
    \end{exerpart}
\end{exer}

\begin{exer} [\exerU] % squaring polynomials]
\label{exer_square_polys}
    Expand and simplify.
    \begin{exerpart}
        \item $(n+1)^2$.
        \item $(2k+1)^2$.
        \item $(m+2)^2-(m-1)^2$
    \end{exerpart}
\end{exer}

\begin{exer}[\exerR] % Polynomials
\label{exer_dyk_polys}
    Do you know \ldots
        \begin{exerpart}
            \item What a polynomial is?
            \item How to multiply polynomials or square a polynomial?
            \item What the abbreviation FOIL stands for?
        \end{exerpart}
\end{exer}

\begin{exer}[\exerE] % last digit of square]  
\label{exer_last_digit_square}
    By \Cref{exer_cases_mod3_mod10}(\ref{part_div_alg_mod10}), we can write a positive integer $a$ as $a=10q+s$ for some integers $q$ and $s$ where $s$ is the last digit of $a$. Note that $s=0,1,2,\ldots,9$.
        \begin{exerpart}
            \item Expand and simplify $a^2$ in terms of $q$ and $s$.
            \item \label{part_last_digit_sq} Show that $a^2-s^2$ is a multiple of 10. It follows that $a^2$ \fn{mod} 10 $=s^2$ \fn{mod} 10 and therefore the last digit of $a^2$ is equal to the last digit of $s^2$.
            \item \label{part_digit_sq} What are the possible last digits of $s^2$?  Hint: Recall that $s$ is a digit, so you can check each value of $s=0,1,2,\ldots,9$.
            \item What are the possible last digits of any integer $a^2$?  Hint: Use the results of (\ref{part_last_digit_sq}) and (\ref{part_digit_sq}).
        \end{exerpart}
\end{exer}

%%%%%%%%%%%%%%%%%%%%%%%%

\subsubsection*{Exercises for \Cref{sub_pff_cases} \subpffcases}

\begin{exer}[\exerP] % proof by cases: quadratic even
\label{exer_quadratic_even}
   For any integer $n$, copy and edit the proof in \Cref{act_proof_cases_nsqplusn_even} to prove that $$n^2+3n+1$$ is odd for any integer $n$.  
\end{exer}

\begin{exer}[\exerU] % proof by cases: $3 \mid n(n+1)(n-1)$]
\label{exer_proof_cases_3divides} 
    Copy and edit the proof in \Cref{exam_product_3consec_mult3} to prove $$3 \mid n(n+1)(n-1)$$ for any integer $n$.   
\end{exer}

\begin{exer}[\exerR] % Proof format: proof by cases
\label{exer_dyk_pff_cases}
    Do you know \ldots
        \begin{exerpart}
            \item How to write a proof by cases?
            \item What needs to be true about the list of cases in a proof-by-cases?
            \item How we can use the Division Algorithm to create cases for a proof by cases?
        \end{exerpart}
\end{exer}

\begin{exer}[\exerE] % Proof not a multiple of three.
\label{exer_proof_cases_not_mult_three}
~
    Use proof by cases (\Cref{pff_cases}) to prove $$3 \nmid 2n^2 +n+1$$ for any integer $n$. Hint: Use \Cref{exer_cases_mod3_mod10} (\ref{part_div_alg_mod3}) to create three cases.  Notice that we are proving something is not a multiple of three.
\end{exer}

\begin{exer}[\exerE] % proof $\lor$ by cases:  $n^2$ ~\fn{mod}~ $4 = 0$ or 1]
\label{exer_proof_cases_square_mod4}
    Use cases (\Cref{pff_cases}) to prove that $n^2$ \fn{mod} $4 = 0$ or $n^2$ \fn{mod} $4 = 1$ for any integer $n$.  
    Hint: Use cases based on whether $n$ is even or $n$ is odd.  Note that in one case we get $n^2$ \fn{mod} $4 = 0$ and in the other case we get $n^2$ \fn{mod} $4 = 1$.  Officially, we are using a modified version of proof by cases that we revisit in \Cref{sec_logic_connectives}.
\end{exer}

\begin{exer}[\exerE] % proof $\lor$ by cases:   $n^2$ ~\fn{mod}~ $3 = 0$ or 1]
\label{exer_proof_cases_square_mod3}
    Use cases (\Cref{pff_cases}) to prove that $n^2$ \fn{mod} $3 = 0$ or $n^2$ \fn{mod} $3 = 1$ for any integer $n$. Hint: Use \Cref{exer_cases_mod3_mod10} (\ref{part_div_alg_mod3}) and consider three cases.  Note that in one case we get $n^2$ \fn{mod} $3 = 0$ and in the other two cases we get $n^2$ \fn{mod} $3 = 1$.  Officially, we are using a modified version of proof by cases that we revisit in \Cref{sec_logic_connectives}.
\end{exer}

%%%%%%%%%%%%%%%%%%%%%%%%%%%
\subsection{\answers}
%%%%%%%%%%%%%%%%%%%%%%%%%%%

%%%%%%%%%%%%%%%%%%%%%%%%%%%%%%%%
\subsubsection{\answers for \Cref{sub_division_algorithm} \subdivalg}

\subsubsection*{\Cref{exer_find_quotient_remainder} \answers}
\begin{apart}
    \item $100 \fn{ div } 3 = 33$ and $100 \fn{ mod }3 = 1$
    \item 2, 6
    \item Hint: Imagine 7 children.
\end{apart}

\subsubsection*{\Cref{exer_examples_mod} \answers}
\begin{apart}
    \item Hint: Check that your answer is between 100 and 200 and odd.
    \item Hint: Check that your answer is one more than a multiple of three.
    \item Hint: for example, $7 \times 20 = 140$, and so 141 is a possible example.
    \item Hint: Any example is of the form $1d1$ where $d$ is a digit.
\end{apart}

\subsubsection*{\Cref{exer_mod10_practice} \answers}
\begin{apart}
    \item Hint: The answers in some order are $0,1,3,7,9$.
    \item Hint: One way to build an example is to pick $a$ and use $b=a+5$.
    \item Hint: Neither $n$ nor $m$ should be a multiple of 10.  Try picking an even integer for $n$ and an odd multiple of five for $m$.
\end{apart}

\subsubsection*{\Cref{exer_modn} \answers}
\begin{apart}
    \item Hint: The answers in some order are 0 and $n$.
    \item Hint: The answers in some order are 0 and 1.
    \item Hint: The answers in some order are 0 and $a$.
\end{apart}

\subsubsection*{\Cref{exer_cases_mod3_mod10} \answers}
\begin{apart}
    \item Hint: By the Division Algorithm, $a=3q+r$ where $0 \le r < 3$.  What are the possible values of the integer $r$?
    \item Hint: Explain why $0 \le s < 10$ and so $s$ is the last digit of $n$.
\end{apart}

\subsubsection*{\Cref{exer_isbn} \answers}
\begin{apart}
    \item Hint: Calculate $(9 + 3 \cdot 7 +8 + 3 \cdot 0 + \ldots ) \fn{ mod } 10.$  You should get zero.
    \item $x=4$
\end{apart}

%%%%%%%%%%%%%%%%%%%%%%%
\subsubsection{\answers for \Cref{sub_polys} \subpolys}

\subsubsection*{\Cref{exer_expanding_polys} \answers}

\begin{apart}
    \item $2n^3+n$
    \item Hint: The coefficient of $k$ in the answer is 4.
    \item Hint: The constant term of the answer is $-1$.
\end{apart}

\subsubsection*{\Cref{exer_square_polys} \answers}
\begin{apart}
    \item Hint: The coefficient of $n$ in the answer is 2.
    \item Hint: In the answer $k^2$ and $k$ have the same coefficient.
    \item $6m+3$
\end{apart}

\subsubsection*{\Cref{exer_cases_mod3_mod10} \answers}
\begin{apart}
    \item Hint: $a^2=\underline{\hspace{0.5in}}q^2 + \underline{\hspace{0.5in}}qs + \underline{\hspace{0.5in}}s^2$
    \item Hint: Expand and simplify $a^2-s^2$ and factor out 10.
    \item Hint: Copy and complete \Cref{table_last_digit_sq}.
        \begin{table}[hbtp]
            \centering
            \begin{tabular}{c |c|c|c |c|c|c |c|c|c |c}
                 $s$ & 0 & 1 & 2 & 3 & 4 & 5 & 6 & 7 & 8 & 9 \\ \hline
                 $s^2$ &&&&& 16 &&&&  \\ \hline
                 Last digit of $s^2$ &&&&& 6 &&&&  \\
            \end{tabular}
            \caption{The last digit of squares}
            \label{table_last_digit_sq}
        \end{table}
    \item Hint: Use the hint.
\end{apart}

%%%%%%%%%%%%%%%%%%%%%%%
\subsubsection{\answers for \Cref{sub_pff_cases} \subpffcases}

\subsubsection*{\Cref{exer_quadratic_even} \answers}
    Hint: In the case where $n=2k+1$, expand and simplify $(2k+1)^2+3(2k+1)+1=4k^2+10k+5 = 2(\underline{\hspace{0.5in}})+1$.  Then, fill in the parentheses.  If you need to, set it equal to $2m+1$ and solve for $m$.
    
\subsubsection*{\Cref{exer_proof_cases_3divides} \answers}
    Hint: As in the proof in \Cref{exam_product_3consec_mult3}, keep the product in factored form and find a way to factor out a three from one of the factors.
    
\subsubsection*{\Cref{exer_proof_cases_square_mod4} \answers}
    Hint: Copy and edit the proof in \Cref{act_proof_cases_nsqplusn_even} (\ref{part_proof_cases_read}) and follow the hint.
    
\subsubsection*{\Cref{exer_proof_cases_square_mod3} \answers}
    Hint: Copy and edit the proof in \Cref{exam_product_3consec_mult3} and follow the hint.

